% Part: first-order-logic
% Chapter: beyond
% Section: higher-order-logic

\documentclass[../../../include/open-logic-section]{subfiles}

\begin{document}

\olfileid[hi]{fol}{byd}{hol}

\olsection{उच्चतर-कोटि तर्कशास्त्र}

प्रथम-कोटि तर्कशास्त्र से द्वितीय-कोटि तर्कशास्त्र की ओर जाने पर हम
औपचारिक भाषा के भीतर प्रथम-कोटि !!{domain} की वस्तुओं के समुच्चयों के विषय
में बात कर सके। यहीं क्यों रुकें? उदाहरणतः तृतीय-कोटि तर्कशास्त्र से हम
वस्तुओं के समुच्चयों के समुच्चयों, अथवा शायद ऐसे समुच्चयों से भी व्यवहार कर
सकेंगे जिनमें वस्तुएँ और वस्तुओं के समुच्चय दोनों हों। चतुर्थ-कोटि
तर्कशास्त्र हमें उस प्रकार की वस्तुओं के समुच्चयों की बात करने देगा। जैसा
आपने अनुमान लगाया होगा, इस विचार को मनचाही बार दोहराया जा सकता है।

व्यवहार में उच्चतर-कोटि तर्कशास्त्र को संबंधों के बजाय प्रायः फलनों के पदों
में !!{formula}बद्ध किया जाता है। (स्वाभाविक अभिज्ञानों को मान लें तो यह
भेद सारहीन है।) कुछ मूल ``वर्ग'' $A$, $B$, $C$,~\dots दिए हों (जिन्हें अब
हम ``प्रकार'' कहेंगे), तो निम्न अनुबंध द्वारा नए प्रकार बना सकते हैं:
\begin{quote}
यदि $\sigma$ और $\tau$ परिमित प्रकार हैं, तो $\sigma \to \tau$ भी परिमित
प्रकार है।
\end{quote}
प्रकारों को ऐसी वाक्यविन्यासिक ``पर्चियाँ'' समझिए जो हमारे !!{domain} में
वांछित वस्तुओं का वर्गीकरण करती हैं; $\sigma \to \tau$ उन वस्तुओं का वर्णन
करता है जो प्रकार~$\sigma$ की वस्तुओं को प्रकार~$\tau$ की वस्तुओं में ले
जाने वाले फलन हैं। उदाहरणतः हम ``सत्य'' और ``असत्य'' सत्यता-मानों का प्रकार
$\Omega$, और प्राकृतिक संख्याओं का प्रकार $\Nat$ रखना चाह सकते हैं। तब
$\Nat \to
\Omega$ प्रकार की वस्तुओं को एकपदी संबंध, अथवा $\Nat$ के
उपसमुच्चय मान सकते हैं; $\Nat \to \Nat$ प्रकार की वस्तुएँ प्राकृतिक
संख्याओं से प्राकृतिक संख्याओं में जाने वाले फलन हैं; और $(\Nat \to \Nat) \to \Nat$
प्रकार की वस्तुएँ ``फलनक'' हैं, अर्थात् ऐसे उच्चतर-प्रकार फलन जो
फलनों को संख्याओं में ले जाते हैं।

द्वितीय-कोटि तर्कशास्त्र की तरह उच्चतर-कोटि तर्कशास्त्र को एक प्रकार का
बहु-वर्गीय तर्कशास्त्र माना जा सकता है, जिसमें विचाराधीन वस्तु के प्रत्येक
प्रकार के लिए एक वर्ग होता है। किंतु सामान्यतः उच्चतर-प्रकार तर्कशास्त्र के
वाक्यविन्यास को आधार से ही परिभाषित करना अधिक स्पष्ट है। उदाहरणतः परिमित
प्रकारों का समुच्चय निम्न प्रकार आगमनिक रूप से परिभाषित कर सकते हैं:
\begin{enumerate}
\item $\Nat$ एक परिमित प्रकार है।
\item यदि $\sigma$ और $\tau$ परिमित प्रकार हैं, तो $\sigma
  \to \tau$ भी
  परिमित प्रकार है।
\item यदि $\sigma$ और $\tau$ परिमित प्रकार हैं, तो $\sigma \times
  \tau$ भी
  परिमित प्रकार है।
\end{enumerate}
सहज अर्थ में $\Nat$ प्राकृतिक संख्याओं का प्रकार, $\sigma \to \tau$ प्रकार
$\sigma$ से प्रकार $\tau$ में जाने वाले फलनों का प्रकार, और $\sigma \times \tau$
वस्तु-युग्मों का प्रकार निरूपित करता है जिनमें एक वस्तु $\sigma$ से
और एक $\tau$ से है। तब पदों का समुच्चय निम्न प्रकार आगमनिक रूप से परिभाषित
कर सकते हैं:
\begin{enumerate}
\item प्रत्येक प्रकार $\sigma$ के लिए चरों $x$, $y$, $z$, \dots का भंडार है;
  ये चर प्रकार $\sigma$ के हैं।
\item $\Obj 0$, प्रकार $\Nat$ का पद है।
\item $\Obj S$ (परवर्ती) प्रकार $\Nat \to \Nat$ का पद है।
\item यदि $s$ प्रकार $\sigma$ का पद और $t$ प्रकार $\Nat \to (\sigma \to \sigma)$
  का पद है, तो $\Obj{R}_{st}$ प्रकार $\Nat \to \sigma$ का पद है।
\item यदि $s$ प्रकार $\tau \to \sigma$ का और $t$ प्रकार~$\tau$ का पद है,
  तो $s(t)$ प्रकार $\sigma$ का पद है।
\item यदि $s$ प्रकार~$\sigma$ का पद और $x$ प्रकार~$\tau$ का चर है, तो
  $\lambd[x][s]$ प्रकार $\tau \to \sigma$ का पद है।
\item यदि $s$ प्रकार~$\sigma$ और $t$ प्रकार~$\tau$ का पद है, तो
  $\tuple{s, t}$ प्रकार $\sigma \times
  \tau$ का पद है।
\item यदि $s$ प्रकार~$\sigma \times \tau$ का पद है, तो $p_1(s)$ प्रकार
  $\sigma$ का और $p_2(s)$ प्रकार~$\tau$ का पद है।
\end{enumerate}
सहज अर्थ में $\Obj{R}_{st}$ निम्न पुनरावर्ती ढंग से परिभाषित फलन निरूपित करता है:
\begin{align*}
\Obj{R}_{st}(0) & = s \\
\Obj{R}_{st}(x+1) & = t(x, R_{st}(x)),
\end{align*}
$\tuple{s, t}$ उस युग्म को निरूपित करता है जिसका पहला घटक~$s$ और दूसरा
घटक~$t$ है; तथा $p_1(s)$ और~$p_2(s)$, $s$ के क्रमशः पहले और दूसरे अवयव
(``प्रक्षेप'') निरूपित करते हैं। अंततः $\lambd[x][s]$ उस फलन~$f$ को निरूपित
करता है जिसकी परिभाषा है
\[
f(x) = s
\]
प्रत्येक~$x$ के लिए जो प्रकार~$\sigma$ का हो; इसलिए मद (6) हमें परिग्रहण का
एक रूप देती है, जिससे हम पदों द्वारा फलन
परिभाषित कर सकते हैं। !!^{formula}s का निर्माण एक ही प्रकार के पदों के बीच
!!{identity}-कथनों $\eq[s][t]$, सामान्य प्रतिज्ञप्तिक संयोजकों और
उच्चतर-प्रकार परिमाणीकरण से होता है। तब तंत्र के अभिगृहीत ऊपर परिभाषित पदों
को नियंत्रित करने वाले मूल समीकरण, तथा परिमाणकों और !!{identity} वाले
तर्कशास्त्र के सामान्य नियम लिए जा सकते हैं।

यदि परिमित-प्रकार तंत्र में सत्यता-मानों का प्रकार $\Omega$ जोड़ें, तो उसके
प्रयोग को नियंत्रित करने वाले अभिगृहीत भी जोड़ने होंगे। वस्तुतः थोड़ी चतुराई
से जटिल !!{formula}s को पूरी तरह हटाकर उनकी जगह प्रकार $\Omega$ के पद रखे जा
सकते हैं!{} फिर प्रमाण-तंत्र को उसी के अनुसार बदला जा सकता है। परिणाम मूलतः
1930 के दशक में अलोंज़ो चर्च द्वारा प्रस्तुत \emph{प्रकारों का सरल सिद्धांत}
है।

द्वितीय-कोटि तर्कशास्त्र की तरह उच्चतर-प्रकार अर्थविज्ञान के भी अलग-अलग रूप
हैं जिनका हम प्रयोग करना चाह सकते हैं। पूर्ण रूप में प्रकार $\sigma \to \tau$
के चर, प्रकार~$\sigma$ की वस्तुओं से प्रकार~$\tau$ की वस्तुओं में
जाने वाले \emph{सभी} फलनों के समुच्चय पर विचरते हैं। अपेक्षानुसार यह
अर्थविज्ञान इतना प्रबल है कि कोई पूर्ण, प्रभावी !!{derivation}-तंत्र स्वीकार
नहीं करता। किंतु दुर्बलतर अर्थविज्ञान लिया जा सकता है, जिसमें
!!a{structure} में अवयवों के समुच्चय $T_\tau$ प्रत्येक प्रकार $\tau$ के लिए
होते हैं, और अनुप्रयोग, प्रक्षेप आदि की उपयुक्त संक्रियाएँ भी होती हैं। विवरण ठीक से
पूरा करने पर ऊपर वर्णित प्रकार के !!{derivation}-तंत्रों के लिए पूर्णता
प्रमेय प्राप्त किए जा सकते हैं।

उच्चतर-प्रकार तर्कशास्त्र आकर्षक है क्योंकि वह ऐसा ढाँचा देता है जिसमें
गणित के बड़े भाग को स्वाभाविक ढंग से अंतःस्थापित कर सकते हैं: $\Nat$ से
आरंभ करके वास्तविक संख्याएँ, संतत फलन आदि परिभाषित किए जा सकते हैं। वह
अंतःप्रज्ञावादी तर्कशास्त्र के संदर्भ में विशेष रूप से आकर्षक भी है, क्योंकि
प्रकारों की स्पष्ट ``रचनात्मक'' व्याख्याएँ हैं। वस्तुतः उच्चतर-प्रकार
अर्थविज्ञान के ऐसे रचनात्मक रूप (चिरसम्मत के बजाय अंतःप्रज्ञावादी
तर्कशास्त्र पर आधारित) विकसित किए जा सकते हैं जो इन रचनात्मक व्याख्याओं को
बहुत अच्छी तरह स्पष्ट करते हैं और अनेक दृष्टियों से अपने चिरसम्मत समकक्षों
से अधिक रोचक हैं।

\end{document}
